\documentclass[12pt,a4paper]{article}

\usepackage[utf8]{inputenc}
\usepackage[T1]{fontenc}
\usepackage[portuguese]{babel}

\usepackage{listings}
\usepackage{xcolor}
\usepackage{amsmath}
\usepackage{amssymb}
\usepackage{amsfonts}
\usepackage{graphicx}

\lstset{
    basicstyle=\ttfamily\small,
    keywordstyle=\color{blue}\bfseries, 
    commentstyle=\color{green!70!black},
    stringstyle=\color{red}, 
    numbers=left, 
    numberstyle=\tiny\color{gray}, 
    frame=single, 
    breaklines=true, 
    showstringspaces=false 
}

\usepackage[margin=2.5cm]{geometry}

\begin{document}

\begin{center}
    {\LARGE MC521 - Relatório V} \\ [2,5em]
    {\large Júlio da Silva Telles} \\ [1,0em]
    {\large 20 de Junho de 2026} \\ [1,5em]
\end{center}

\vspace{1cm} 

\section{C - Orac and Models (CodeForces - 1350B)}
O problema consiste em encontrar o tamanho máximo de uma sequência de índices $i_1 < i_2 < \dots < i_k$, tal que cada índice na sequência seja um múltiplo do índice anterior (ou seja, $i_x$ divide $i_{x+1}$) e os valores correspondentes no vetor original sejam estritamente crescentes ($v[i_x] < v[i_{x+1}]$).

\subsection*{1.1 Solução}
Podemos resolver este problema utilizando Programação Dinâmica (DP) de forma iterativa, com uma complexidade de tempo de $\mathcal{O}(N \log N)$ graças à soma da série harmônica ao iterar pelos múltiplos. O algoritmo segue os seguintes passos:

\begin{enumerate}
    \item \textbf{Inicialização do Estado:} Inicializamos um vetor de estados da DP (\texttt{s}) com o valor base 1 para todos os elementos a partir da metade do array, pois o tamanho mínimo de qualquer sequência válida é o próprio elemento sozinho.
    \item \textbf{Iteração e Transição (Bottom-Up):} Percorremos os índices do array de trás para frente, de $N/2$ até 1. Para cada índice, verificamos todos os seus múltiplos possíveis (multiplicando o índice por $j = 2, 3, \dots$).
    \item \textbf{Validação e Atualização:} Para cada múltiplo válido, verificamos se o valor no array original satisfaz a condição de crescimento estrito. Se for menor, testamos se incluir este múltiplo melhora a nossa sequência local (\texttt{s[i] < 1 + s[i * j]}). Se melhorar, atualizamos o valor de \texttt{s[i]} e guardamos o tamanho máximo global encontrado para impressão no final de cada caso de teste.
\end{enumerate}

\section{G - Partial Teacher (CodeForces-67A)}
O problema exige a atribuição de valores a um conjunto de $V$ vértices de forma a respeitar condições de desigualdade direcionais ("L" para menor, "R" para maior) ou igualdade ("=") fornecidas em formato de string. O objetivo é calcular o valor acumulado ou o tamanho do maior caminho direcionado para cada nó, respeitando essas relações.

\subsection*{2.1 Solução}
Podemos modelar o problema como a busca do caminho mais longo em um Grafo Direcionado Acíclico (DAG) utilizando Busca em Profundidade (DFS) com \textit{Memoization} (Programação Dinâmica em grafos). O algoritmo segue os seguintes passos:

\begin{enumerate}
    \item \textbf{Construção do Grafo:} Lemos a string de relações e interpretamos os caracteres como arestas. Se a relação for "L", criamos uma aresta do nó atual para o próximo (indicando precedência). Se for "R", a direção da aresta é invertida. Se for "=", armazenamos essa equivalência em um vetor auxiliar (\texttt{equals}) em vez de criar uma aresta.
    \item \textbf{Busca em Profundidade (Memoization):} Para cada vértice do grafo, iniciamos uma travessia DFS (\texttt{dpfs}). Durante a recursão, visitamos todos os nós adjacentes e calculamos o tamanho do maior caminho encontrado. O resultado é armazenado no vetor \texttt{sol} para evitar que subproblemas sejam recalculados desnecessariamente.
    \item \textbf{Tratamento de Igualdades e Retorno:} A base da recursão inclui a verificação do vetor \texttt{sol}. Ao terminarmos de calcular o valor de um nó, checamos se ele possui um nó adjacente marcado como igual. Se sim, propagamos o mesmo valor máximo de caminho para o nó equivalente. Por fim, imprimimos o valor processado para todos os vértices.
\end{enumerate}

\end{document}